\begin{afterword}
\summaryhead

The post begins with a complaint the author says he often hears: people who reject his ideas about superintelligence and Friendly AI say he should not be believed because he has no PhD. He argues that this is rarely their real reason. The real rejection happens first and fast, by pattern-matching the ideas to science fiction or a doomsday cult. The PhD is a justification found afterwards. If he had a PhD, the same people would ask for tenure, or a Nobel prize. He cites Eric Drexler, who earned a doctorate for his nanotechnology work without, as far as the author knows, changing his critics' minds.

He generalizes to business and to disagreement between rationalists. Lasting disagreements usually have sources that are hard to state, such as intuitions, professional habits or self-deception. Each side should ask itself ``Is this my true rejection?'' He adds that the question may also be put to the other person, if it is put humbly.

\responsehead

The post turns a personal grievance into a general tool, and the tool works in one direction. Most of the post is spent showing that the author's critics do not mean what they say. The advice to question oneself comes in the third-to-last paragraph and is reversed in the last.

The tool also has an older name. To give acceptable reasons that hide the real cause of a judgment is what psychologists since Ernest Jones in 1908 have called rationalization. Fourteen months earlier, in ``Rationalization,'' the author called that word a ``wrong word'' coined by a ``fool.'' Here he needs the concept, and he renames it instead of crediting it.

As a tool it has the fault the author warned against in ``The Bottom Line.'' He wrote there that his warning was ``not a Fully General Counterargument against conclusions you don't like.'' ``That is not your true rejection'' is exactly such a counterargument. It can be said to anyone, about any objection, and the listener cannot refute it, because on the post's own definition the true rejection is not open to introspection. If the author meets one objection and the critic raises another, the new one shows only that the old one was not the true rejection. Drexler's case shows this in motion: his best-known critic gave technical reasons after the doctorate, and the post does not mention them.

The post never considers that its critics might be right about credentials. Where a listener cannot evaluate the argument, trusting the speaker's record is reasonable, and the author said so himself a year earlier. Superintelligence in 2008 was such a domain. The PhD objection was crude, but it pointed at a real question: who, qualified to judge, had checked these claims? The post answers the crude version and not the question.

The context sharpens this. The post appeared during the author's public debate with Robin Hanson, and two days later he used its question on Hanson, ``in all humility.'' Reading the post as a move in that debate is an inference from the dates. But the post's first use and its later use point the same way: at someone else.

In short: the author rediscovers rationalization, gives it a new name, and uses it to show that his critics' stated reasons do not count, while a Nobel laureate's technical objections to his own example go unmentioned.
\end{afterword}
